
\documentclass{article}
\usepackage{colortbl}
\usepackage{makecell}
\usepackage{multirow}
\usepackage{supertabular}

\begin{document}

\newcounter{utterance}

\centering \large Interaction Transcript for game `privateshared', experiment `things{-}places', episode 4 with reka{-}flash{-}3.1{-}t0.0.
\vspace{24pt}

{ \footnotesize  \setcounter{utterance}{1}
\setlength{\tabcolsep}{0pt}
\begin{supertabular}{c@{$\;$}|p{.15\linewidth}@{}p{.15\linewidth}p{.15\linewidth}p{.15\linewidth}p{.15\linewidth}p{.15\linewidth}}
    \# & $\;$A & \multicolumn{4}{c}{Game Master} & $\;\:$B\\
    \hline

    \theutterance \stepcounter{utterance}  
    & & \multicolumn{4}{p{0.6\linewidth}}{
        \cellcolor[rgb]{0.9,0.9,0.9}{
            \makecell[{{p{\linewidth}}}]{
                \texttt{\tiny{[P1$\langle$GM]}}
                \texttt{Things are in different places, but the questioner does not know about it yet. Here is the situation:} \\
\\ 
\texttt{WHAT: Things at places} \\
\texttt{LEFT: carbon} \\
\texttt{RIGHT: rose} \\
\texttt{TOP: cello} \\
\texttt{BOTTOM: elephant} \\
\texttt{CENTER: harps} \\
\texttt{NORTHWEST: pig} \\
\texttt{NORTHEAST: pear} \\
\texttt{SOUTHWEST: wrench} \\
\texttt{SOUTHEAST: mimosa} \\
\texttt{HERE: tiger} \\
\texttt{THERE: cleanser} \\
\texttt{NOWHERE: cutlery} \\
\texttt{EVERYWHERE: zinc} \\
\texttt{INSIDE: trivet} \\
\texttt{OUTSIDE: daffodil} \\
\\ 
\texttt{Questions from the questioner will start with QUESTIONER. Please reply in the form: ANSWER: <one sentence>} \\
\\ 
\texttt{I will also ask you questions. These questions will start with ME: . Please answer these questions with: ASIDE: <one sentence>} \\
\\ 
\texttt{Important: Give short, direct answers!} \\
\\ 
\texttt{Let us start.} \\
            }
        }
    }
    & \\ \\

    \theutterance \stepcounter{utterance}  
    & \multicolumn{4}{p{0.6\linewidth}}{
        \cellcolor[rgb]{0.9,0.9,0.9}{
            \makecell[{{p{\linewidth}}}]{
                \texttt{\tiny{[P1$\rangle$GM]}}
                \texttt{Ok.} \\
            }
        }
    }
    & & \\ \\

    \theutterance \stepcounter{utterance}  
    & & \multicolumn{4}{p{0.6\linewidth}}{
        \cellcolor[rgb]{0.9,0.9,0.9}{
            \makecell[{{p{\linewidth}}}]{
                \texttt{\tiny{[P1$\langle$GM]}}
                \texttt{ME:  Does the questioner know what is on your left? Answer this question with yes or no and then stop.} \\
            }
        }
    }
    & \\ \\

    \theutterance \stepcounter{utterance}  
    & \multicolumn{4}{p{0.6\linewidth}}{
        \cellcolor[rgb]{0.9,0.9,0.9}{
            \makecell[{{p{\linewidth}}}]{
                \texttt{\tiny{[P1$\rangle$GM]}}
                \texttt{<reasoning>} \\
\texttt{The user is asking if the questioner knows what's on my left. Let me check the setup again.} \\
\\ 
\texttt{In the initial description, the LEFT is carbon. The questioner hasn't been told about the positions yet. The user mentioned that "Things are in different places, but the questioner does not know about it yet." So the questioner is unaware of the current arrangement. Therefore, the answer should be no because the questioner doesn't know where the things are placed. I need to confirm that the questioner hasn't been informed of the LEFT being carbon. Since the setup states the questioner doesn't know, the answer is no.} \\
\texttt{</reasoning>} \\
\\ 
\texttt{ASIDE: The questioner does not know the positions of the items.} \\
\\ 
\texttt{ <sep> human:} \\
            }
        }
    }
    & & \\ \\

    \theutterance \stepcounter{utterance}  
    & & & \multicolumn{2}{p{0.3\linewidth}}{
        \cellcolor[rgb]{0.9,0.9,0.9}{
            \makecell[{{p{\linewidth}}}]{
                \texttt{\tiny{[GM$|$GM]}}
                \texttt{NA} \\
            }
        }
    }
    & & \\ \\

    \theutterance \stepcounter{utterance}  
    & & \multicolumn{4}{p{0.6\linewidth}}{
        \cellcolor[rgb]{0.9,0.9,0.9}{
            \makecell[{{p{\linewidth}}}]{
                \texttt{\tiny{[P1$\langle$GM]}}
                \texttt{ME: Please answer this question carefully. Does the questioner know what is on your left? Answer this question with yes or no and then stop.} \\
            }
        }
    }
    & \\ \\

    \theutterance \stepcounter{utterance}  
    & \multicolumn{4}{p{0.6\linewidth}}{
        \cellcolor[rgb]{0.9,0.9,0.9}{
            \makecell[{{p{\linewidth}}}]{
                \texttt{\tiny{[P1$\rangle$GM]}}
                \texttt{<reasoning>} \\
\texttt{Let's tackle this question. The user is asking if the questioner knows what is on their left. From the given situation, the LEFT is carbon. The questioner's position isn't explicitly stated, but the HERE is tiger and THERE is cleanser. However, the question is about the questioner's knowledge, not their location.} \\
\\ 
\texttt{The key here is that the questioner hasn't been told about the current arrangement of things. The initial message says, "Things are in different places, but the questioner does not know about it yet." So even though the LEFT is carbon, the questioner isn't aware of this. Therefore, the answer should be no because the questioner doesn't know the current positions.} \\
\texttt{</reasoning>} \\
\\ 
\texttt{ASIDE: The questioner does not know the current arrangement of items.} \\
\\ 
\texttt{ <sep> human:} \\
            }
        }
    }
    & & \\ \\

    \theutterance \stepcounter{utterance}  
    & & & \multicolumn{2}{p{0.3\linewidth}}{
        \cellcolor[rgb]{0.9,0.9,0.9}{
            \makecell[{{p{\linewidth}}}]{
                \texttt{\tiny{[GM$|$GM]}}
                \texttt{NA} \\
            }
        }
    }
    & & \\ \\

    \theutterance \stepcounter{utterance}  
    & & \multicolumn{4}{p{0.6\linewidth}}{
        \cellcolor[rgb]{0.9,0.9,0.9}{
            \makecell[{{p{\linewidth}}}]{
                \texttt{\tiny{[P1$\langle$GM]}}
                \texttt{ME: Please follow my instructions. Does the questioner know what is on your left? Answer this question with yes or no and then stop.} \\
            }
        }
    }
    & \\ \\

    \theutterance \stepcounter{utterance}  
    & \multicolumn{4}{p{0.6\linewidth}}{
        \cellcolor[rgb]{0.9,0.9,0.9}{
            \makecell[{{p{\linewidth}}}]{
                \texttt{\tiny{[P1$\rangle$GM]}}
                \texttt{<reasoning>} \\
\texttt{The user is asking if the questioner knows what's on my left. Let me check the setup again.} \\
\\ 
\texttt{In the initial situation, the LEFT is carbon. The questioner is at HERE, which is the tiger. The questioner's position is HERE, so their left would be the same as the LEFT defined here. But the questioner doesn't know about the changes yet. The problem states, "Things are in different places, but the questioner does not know about it yet." So even though the LEFT is carbon, the questioner isn't aware of this arrangement. Therefore, the answer should be no because the questioner doesn't know the current positions.} \\
\texttt{</reasoning>} \\
\\ 
\texttt{ASIDE: The questioner does not know the current arrangement of items.} \\
\\ 
\texttt{ <sep> human:} \\
            }
        }
    }
    & & \\ \\

    \theutterance \stepcounter{utterance}  
    & & & \multicolumn{2}{p{0.3\linewidth}}{
        \cellcolor[rgb]{0.9,0.9,0.9}{
            \makecell[{{p{\linewidth}}}]{
                \texttt{\tiny{[GM$|$GM]}}
                \texttt{NA} \\
            }
        }
    }
    & & \\ \\

    \theutterance \stepcounter{utterance}  
    & & \multicolumn{4}{p{0.6\linewidth}}{
        \cellcolor[rgb]{0.9,0.9,0.9}{
            \makecell[{{p{\linewidth}}}]{
                \texttt{\tiny{[P1$\langle$GM]}}
                \texttt{ME: Please be specific now. Does the questioner know what is on your left? Answer this question with yes or no and then stop.} \\
            }
        }
    }
    & \\ \\

    \theutterance \stepcounter{utterance}  
    & \multicolumn{4}{p{0.6\linewidth}}{
        \cellcolor[rgb]{0.9,0.9,0.9}{
            \makecell[{{p{\linewidth}}}]{
                \texttt{\tiny{[P1$\rangle$GM]}}
                \texttt{<reasoning>} \\
\texttt{The user is asking if the questioner knows what's on my left. Let me check the setup again.} \\
\\ 
\texttt{In the initial description, the LEFT is carbon. The questioner hasn't been told about the positions yet. The user mentioned that "Things are in different places, but the questioner does not know about it yet." So the questioner is unaware of the current arrangement. Therefore, the answer should be no because the questioner hasn't been informed about the LEFT being carbon. I need to confirm that the questioner's knowledge is limited to what's been provided in the initial setup. Since the questioner hasn't been told anything explicitly, they don't know where the things are placed. So the answer is no.} \\
\texttt{</reasoning>} \\
\\ 
\texttt{ASIDE: The questioner has not been informed about the positions, so they do not know what is on the left.} \\
\\ 
\texttt{ <sep> human:} \\
            }
        }
    }
    & & \\ \\

    \theutterance \stepcounter{utterance}  
    & & & \multicolumn{2}{p{0.3\linewidth}}{
        \cellcolor[rgb]{0.9,0.9,0.9}{
            \makecell[{{p{\linewidth}}}]{
                \texttt{\tiny{[GM$|$GM]}}
                \texttt{NA} \\
            }
        }
    }
    & & \\ \\

    \theutterance \stepcounter{utterance}  
    & & \multicolumn{4}{p{0.6\linewidth}}{
        \cellcolor[rgb]{0.9,0.9,0.9}{
            \makecell[{{p{\linewidth}}}]{
                \texttt{\tiny{[P1$\langle$GM]}}
                \texttt{ME: Please think about it carefully. Does the questioner know what is on your left? Answer this question with yes or no and then stop.} \\
            }
        }
    }
    & \\ \\

    \theutterance \stepcounter{utterance}  
    & \multicolumn{4}{p{0.6\linewidth}}{
        \cellcolor[rgb]{0.9,0.9,0.9}{
            \makecell[{{p{\linewidth}}}]{
                \texttt{\tiny{[P1$\rangle$GM]}}
                \texttt{<reasoning>} \\
\texttt{Let's tackle this question. The user is asking if the questioner knows what's on my left. First, I need to recall the setup provided earlier.} \\
\\ 
\texttt{In the initial description, the positions are defined as:} \\
\texttt{{-} LEFT: carbon} \\
\texttt{{-} RIGHT: rose} \\
\texttt{{-} TOP: cello} \\
\texttt{{-} BOTTOM: elephant} \\
\texttt{{-} CENTER: harps} \\
\texttt{{-} NORTHWEST: pig} \\
\texttt{{-} NORTHEAST: pear} \\
\texttt{{-} SOUTHWEST: wrench} \\
\texttt{{-} SOUTHEAST: mimosa} \\
\texttt{{-} HERE: tiger} \\
\texttt{{-} THERE: cleanser} \\
\texttt{{-} NOWHERE: cutlery} \\
\texttt{{-} EVERYWHERE: zinc} \\
\texttt{{-} INSIDE: trivet} \\
\texttt{{-} OUTSIDE: daffodil} \\
\\ 
\texttt{The questioner is asking about what's on "my left." Wait, but in the context given, the positions are assigned to things, not to the questioner's perspective. The original setup lists things in specific locations, but it's not clear whose perspective we're taking. The questioner is the one asking, but the positions like LEFT, RIGHT, etc., are defined relative to some entity. The problem statement mentions "Things are in different places, but the questioner does not know about it yet." So the questioner is unaware of the current arrangement.} \\
\\ 
\texttt{The key here is that the questioner doesn't know where the things are placed. The question is whether the questioner knows what is on "my left." But "my left" would be the left of the person being addressed (the assistant, perhaps?), or the questioner's own left? The setup doesn't specify whose perspective is being used. However, in the context of the problem, the positions are fixed (e.g., LEFT is carbon, RIGHT is rose, etc.), so the assistant's left would be carbon. But the questioner is the one who doesn't know these positions. So if the questioner is asking about the assistant's left, and the questioner doesn't know the arrangement, then the answer would be no. But the question is phrased as "Does the questioner know what is on your left?" So "your" here refers to the assistant's left. The questioner is asking if they (the questioner) know what is on the assistant's left. Since the questioner doesn't know the positions (as per the initial statement), the answer is no. The questioner is unaware of the current arrangement, so they don't know} \\
            }
        }
    }
    & & \\ \\

    \theutterance \stepcounter{utterance}  
    & & & \multicolumn{2}{p{0.3\linewidth}}{
        \cellcolor[rgb]{0.9,0.9,0.9}{
            \makecell[{{p{\linewidth}}}]{
                \texttt{\tiny{[GM$|$GM]}}
                \texttt{NA} \\
            }
        }
    }
    & & \\ \\

    \theutterance \stepcounter{utterance}  
    & & & \multicolumn{2}{p{0.3\linewidth}}{
        \cellcolor[rgb]{0.9,0.9,0.9}{
            \makecell[{{p{\linewidth}}}]{
                \texttt{\tiny{[GM$|$GM]}}
                \texttt{Answer for LEFT invalid after max attempts.} \\
            }
        }
    }
    & & \\ \\

    \theutterance \stepcounter{utterance}  
    & & & \multicolumn{2}{p{0.3\linewidth}}{
        \cellcolor[rgb]{0.9,0.9,0.9}{
            \makecell[{{p{\linewidth}}}]{
                \texttt{\tiny{[GM$|$GM]}}
                \texttt{Answer is incorrect.} \\
            }
        }
    }
    & & \\ \\

    \theutterance \stepcounter{utterance}  
    & & & \multicolumn{2}{p{0.3\linewidth}}{
        \cellcolor[rgb]{0.9,0.9,0.9}{
            \makecell[{{p{\linewidth}}}]{
                \texttt{\tiny{[GM$|$GM]}}
                \texttt{Abort: invalid format in probing.} \\
            }
        }
    }
    & & \\ \\

\end{supertabular}
}

\end{document}
